\documentclass[11pt]{article}

\usepackage[margin=1in]{geometry}
\usepackage{amsmath,amssymb}
\usepackage{booktabs}
\usepackage{array}
\usepackage{xcolor}
\usepackage[hidelinks]{hyperref}
\usepackage{enumitem}
\usepackage{microtype}
\usepackage[htt]{hyphenat}
\usepackage{seqsplit}
\usepackage[ruled,vlined,linesnumbered]{algorithm2e}

\definecolor{unityblue}{HTML}{2C5F9E}
\definecolor{pyorange}{HTML}{C9560B}
\definecolor{okgreen}{HTML}{1B7A3D}
\newcommand{\code}[1]{{\texttt{\small #1}}}
\newcommand{\unity}[1]{{\color{unityblue}\texttt{\small #1}}}
\newcommand{\py}[1]{{\color{pyorange}\texttt{\small #1}}}
\newcommand{\clip}{\operatorname{clip}_{[0,1]}}
\setlength{\emergencystretch}{3em}

\title{\textbf{The \texttt{rules-based-v2} Reference Policy\\ in the ARC Disaster-Response Environment}}
\author{ARC\_Game RL --- technical reference}
\date{\today}

\begin{document}
\maketitle

\begin{abstract}
\noindent \py{rules-based-v2} (\py{improved\_rules\_based\_decision} in \py{benchmark\_models.py})
is a non-learning, API-free heuristic that repairs the four documented flaws of the baseline
\py{rules-based} policy: it builds and hires \emph{multiple} units per turn, invests in capacity
\emph{ahead} of demand (bounded by what it can staff), chooses build sites \emph{geographically},
and \emph{deploys} its budget rather than hoarding a fixed reserve. Its task choices use a value
function with the motel's recurring cost \emph{priced in}, so relocations route into free shelters
automatically. In an 8-episode smoke it scores $\approx 3.0$ vs the baseline's $2.5$ (full food,
$93\%$ lodging, casework $0.43$). Reward terms are defined in the companion \emph{Reward Scoring}
document.
\end{abstract}

\tableofcontents
\vspace{1em}

% =====================================================================
\section{Hardcoded Limits and Constants}
% =====================================================================

The policy has two kinds of constants: \textbf{policy hyperparameters} it chooses, and
\textbf{environment limits} it must respect (these are enforced by the simulator / action
enumerator, not by the policy). Table~\ref{tab:consts} lists both.

\renewcommand{\arraystretch}{1.25}
\begin{table}[h]
\centering
\footnotesize
\begin{tabular}{>{\raggedright\arraybackslash}p{3.5cm} l >{\raggedright\arraybackslash}p{8.0cm}}
\toprule
\textbf{Symbol (code)} & \textbf{Value} & \textbf{Role} \\
\midrule
\multicolumn{3}{l}{\textit{Policy hyperparameters (set in \py{benchmark\_models.py})}}\\
\midrule
$B_{\text{turn}}$ (\py{\_V2\_BUILD\_PER\_TURN}) & $3$ & Max new buildings appended in one turn. \\
$U_{\max}$ (\py{\_V2\_PIPELINE\_CAP}) & $4$ & Max concurrent \emph{unstaffed} buildings (UnderConstruction / NeedWorker) allowed before building pauses. \\
$\beta$ (\py{\_V2\_OP\_BUFFER}) & $\$3000$ & Cash kept before any discretionary build/hire --- sized to fund a reactive paid food airlift, \emph{not} hoarded. \\
$K_{\text{tgt}}$ (\py{\_POT\_KITCHEN\_TARGET}) & $2$ & Target number of kitchens. \\
$\rho$ (\py{\_ROUNDS\_PER\_DAY}) & $4$ & Rounds per in-game day (motel bills per day). \\
$\sigma$ (\py{\_V2\_SHELTER\_BEDS}) & $10$ & Population capacity per shelter (for sizing). \\
$\delta_{\text{shel}}$ (in-code) & $100$ & Penalty subtracted from a shelter choice that lacks space. \\
\midrule
\multicolumn{3}{l}{\textit{Environment limits the policy respects (enforced by the simulator)}}\\
\midrule
$H_{\text{day}}$ (\py{DAILY\_HIRE\_LIMIT}) & $5$/day & Max workers hired per in-game day; hire actions bundle quantities $1$–$5$. \\
$\omega$ (\py{BUILDING\_REQUIRED\_WORKFORCE}) & $4$ & Workers to operate a building (used to size staffing). \\
$m$ (motel rate) & $\$200$/person/day & Recurring motel charge priced into the choice value. \\
$w_c$ (\py{w\_food\_cost}) & $0.0002$ & Reward cost weight, reused so the choice value matches the reward scale. \\
\bottomrule
\end{tabular}
\caption{Constants in \py{rules-based-v2}. Policy hyperparameters are tunable; environment limits
are fixed by the simulator and merely respected by the policy.}
\label{tab:consts}
\end{table}

% =====================================================================
\section{Decision Pipeline}
% =====================================================================

Each round the policy returns task-choice selections $C$ and a list of game-action indices $X$,
built in five stages:
\begin{enumerate}[leftmargin=1.6em, itemsep=1pt]
\item \textbf{Reactive core} --- inherit the greedy policy's zero-cost worker assignments (staff
buildings that need workers).
\item \textbf{Long-term task choices} --- answer every open task with a choice scored by a value
function that prices in the motel's recurring cost (\S\ref{sec:choice}).
\item \textbf{Staffing-aware capacity target} --- decide \emph{what} to build, bounded by what can
be staffed over the horizon (\S\ref{sec:invest}).
\item \textbf{Geographic, multi-build construction} --- decide \emph{where} and append several
builds, paced by a staffable pipeline and the cash buffer (\S\ref{sec:build}).
\item \textbf{Multi-hire} --- hire the largest affordable batch to close the staffing gap
(\S\ref{sec:hire}).
\end{enumerate}

% =====================================================================
\section{Stage 2: The Long-Term Choice Value Function}
\label{sec:choice}
% =====================================================================

For an active task $t$ with demand indicator $D(t)=\mathbf{1}[\text{taskType}\in\{\text{Demand},\text{Emergency}\}]$,
each choice $c$ exposes a budget impact $b(c)$ and satisfaction impact $s(c)$. Choices are first
\emph{classified by destination} from the choice text $\tau(c)$ (lower-cased), guarding against a
client group's \emph{name} containing ``Motel''/``Shelter'':
\begin{align}
\text{cw}(c) &= [\text{``casework''} \in \tau(c)], \\
\text{mo}(c) &= [\text{``motel''} \in \tau(c)] \wedge \neg\,\text{cw}(c), \\
\text{sh}(c) &= [\text{``shelter''} \in \tau(c)] \wedge \neg\,\text{mo}(c) \wedge \neg\,\text{cw}(c).
\end{align}
With $\mathrm{cost}(c)=\max(0,-b(c))$ and the \emph{acting} predicate
$\alpha(c) = [\mathrm{cost}(c)>0] \vee [s(c)\ge 10]$, the value is
\begin{equation}
v(c) =
\begin{cases}
\dfrac{b(c)}{10^4} + 10^{-2} s(c), & b(c) > 0\ \ \text{(funding)}\\[10pt]
\mathbf{1}[\alpha(c)\wedge D(t)] + 10^{-2}s(c) - w_c\big(\mathrm{cost}(c) + R(c)\big), & b(c) \le 0\ \ \text{(acting/waiting)}
\end{cases}
\label{eq:v2val}
\end{equation}
where the key addition over the myopic greedy value is the \textbf{long-term recurring cost} of the
motel, charged over the remaining days of the episode:
\begin{equation}
R(c) = \begin{cases} m \cdot q(c) \cdot \max\!\big(1,\ \text{roundsLeft}/\rho\big), & \text{mo}(c)\\[4pt] 0, & \text{otherwise,}\end{cases}
\end{equation}
with $m = \$200$/person/day and $q(c)$ the relocation quantity (\code{deliveryQuantity}, default $20$).
Because $w_c\,R(c)$ is large early (many days left) and shrinks to near zero late, shelters
dominate the motel for most of the episode and the motel becomes acceptable only near the end ---
so \emph{relocation routing into shelters emerges from the value function}, with no separate
override. The \emph{acting} predicate ensures free non-acting options (e.g.\ ``request from
kitchens'', which fails without a kitchen) score $0$ and lose to the reliable paid action.

\paragraph{Per-task selection with a shelter-space guard.} For each task the policy takes the
highest-value choice, but first subtracts $\delta_{\text{shel}}=100$ from any shelter choice that
cannot fit the relocation (so it falls back to the motel rather than a deferred shelter delivery
that would fail). Let $g$ be the current free shelter space (sum over \emph{InUse} shelters) and
$p_t$ the affected community's population (default $20$):
\begin{equation}
\hat v(c) = v(c) - \delta_{\text{shel}}\cdot \mathbf{1}\big[\text{sh}(c)\wedge g < p_t\big].
\end{equation}
The best choice $c^\star=\arg\max_c \hat v(c)$ is committed iff $\hat v(c^\star)>0$ \emph{or} it
fulfils a real demand ($\alpha(c^\star)\wedge D(t)$); if a shelter choice is taken, $g$ is decremented
by $p_t$ (Algorithm~\ref{alg:choice}).

\begin{algorithm}[h]
\DontPrintSemicolon
\SetKwInOut{Input}{input}\SetKwInOut{Output}{output}
\Input{active tasks; free shelter space $g$; roundsLeft}
\Output{choices $C$}
$C \leftarrow \emptyset$\;
\ForEach{active task $t$ with choices}{
  $D \leftarrow \mathbf{1}[\text{taskType}\in\{\text{Demand},\text{Emergency}\}]$;\quad $p_t \leftarrow$ affected community pop (default 20)\;
  $c^\star \leftarrow \arg\max_{c}\ \hat v(c)$ \tcp*{Eq.~\eqref{eq:v2val} + shelter-space guard}
  \If{$\hat v(c^\star) > 0$ \textbf{or} $c^\star$ fulfils a demand}{
    add $(t, c^\star)$ to $C$\;
    \If{$c^\star$ is a shelter choice \textbf{and} $g \ge p_t$}{ $g \mathrel{-}= p_t$ }
  }
}
\Return $C$\;
\caption{\textsc{LongTermChoiceSelection}}
\label{alg:choice}
\end{algorithm}

% =====================================================================
\section{Stage 3: Staffing-Aware Capacity Target}
\label{sec:invest}
% =====================================================================

The policy invests \emph{ahead} of demand, but never queues more capacity than it can \emph{staff}:
hiring is capped at $H_{\text{day}}=5$/day and each building needs $\omega=4$ workers, so chasing
the full displaced population would just create unstaffable buildings. Let $W$ be the current total
workers (free $+$ working), and $\text{roundsLeft}$ the rounds remaining; the staffable building
budget is
\begin{align}
d &= \big\lceil \text{roundsLeft}/\rho \big\rceil \quad\text{(days left)}, &
W_{\max} &= W + H_{\text{day}}\cdot d, \\
N_{\max} &= \big\lfloor W_{\max}/\omega \big\rfloor, &
N_{\text{budget}} &= \max\!\big(0,\ N_{\max} - N_{\text{exist}}\big),
\end{align}
where $N_{\text{exist}} = n_{\text{cw}} + n_{\text{kit}} + n_{\text{shel}}$ counts all current
buildings. The shelter shortfall toward the known displaced population $P$ (sum of community
populations) is $n_{\text{need}} = \big\lceil \max(0, P-\text{shelterCap})/\sigma \big\rceil$.
The build wishlist is then assembled \textbf{kitchens before shelters} (food is a full reward point
hard-gated by a kitchen) and truncated to the staffable budget:
\begin{equation}
\begin{aligned}
\textsc{want} &= \big[\textsc{CaseworkSite}\big]_{n_{\text{cw}}<1} \ \Vert\ \big[\textsc{Kitchen}\big]^{(K_{\text{tgt}}-n_{\text{kit}})} \ \Vert\ \big[\textsc{Shelter}\big]^{\,n_{\text{need}}},\\
\textsc{want} &\leftarrow \textsc{want}[\,:N_{\text{budget}}\,] \quad\text{(truncate to the staffable budget).}
\end{aligned}
\end{equation}

% =====================================================================
\section{Stage 4: Geographic, Multi-Build Construction}
\label{sec:build}
% =====================================================================

\paragraph{Where to build.} Each available site $s$ is ranked by distance from its world position
to the nearest community, with a large penalty for flood-blocked routes (so shelters sit close to
the communities they serve and avoid blocked sites):
\begin{equation}
\operatorname{rank}(s) = \min_{c\,\in\,\text{communities}} \lVert \mathrm{pos}(s) - \mathrm{pos}(c)\rVert_2 \;+\; 10^{6}\cdot \mathbf{1}\big[\mathrm{name}(s)\in \text{blockedRoutes}\big].
\end{equation}

\paragraph{How many.} The policy walks \textsc{want} in priority order and appends a build at the
best-ranked unused site for each type, stopping when any guard trips: the per-turn cap
$B_{\text{turn}}=3$, the unstaffed-pipeline cap $U_{\max}=4$ (current unstaffed $+$ queued), or the
cash buffer (a build is only committed if $\mathrm{cost} \le \text{budget} - \beta$). Building is
skipped entirely when fewer than $2$ rounds remain.

% =====================================================================
\section{Stage 5: Multi-Hire}
\label{sec:hire}
% =====================================================================

The staffing gap counts workers needed by currently-understaffed buildings \emph{plus} four per
building queued this turn, minus the free pool $f$:
\begin{equation}
\text{gap} = \max\!\Big(0,\ \textstyle\sum_{b\,\in\,\text{NeedWorker}}\!\big(\text{required}(b)-\text{assigned}(b)\big) \;+\; \omega\cdot|\text{built}| \;-\; f\Big).
\end{equation}
If $\text{gap}>0$, the policy appends the single \unity{hire\_untrained} action with the
\emph{largest} quantity whose cost fits the buffer ($\le \text{budget}-\beta$). The action
enumerator already caps the offered quantity at the remaining daily limit, so one action hires up
to $5$ workers --- the most the environment permits in a day.

% =====================================================================
\section{The Complete Policy}
% =====================================================================

\begin{algorithm}[h]
\DontPrintSemicolon
\SetKwInOut{Input}{input}\SetKwInOut{Output}{output}
\Input{game state; round $r$; total rounds $T$;\quad $\text{roundsLeft}=\max(0,T-r)$}
\Output{choices $C$; action indices $X$}
$X \leftarrow$ greedy worker assignments (reactive core)\;
$C \leftarrow \textsc{LongTermChoiceSelection}(\cdot)$ \tcp*{Algorithm~\ref{alg:choice}}
\BlankLine
compute $P$, shelterCap, $n_{\text{cw}}, n_{\text{kit}}, n_{\text{shel}}$, $W$, budget\;
$N_{\text{budget}} \leftarrow \max(0, \lfloor (W + 5\lceil \text{roundsLeft}/\rho\rceil)/\omega \rfloor - N_{\text{exist}})$\;
$\textsc{want} \leftarrow ([\textsc{Casework}]_{n_{\text{cw}}<1}\Vert[\textsc{Kitchen}]^{K_{\text{tgt}}-n_{\text{kit}}}\Vert[\textsc{Shelter}]^{n_{\text{need}}})[\,:N_{\text{budget}}]$\;
rank available sites by $\operatorname{rank}(s)$\;
\BlankLine
\If{$\mathrm{roundsLeft} \ge 2$}{
  \ForEach{$\text{btype} \in \textsc{want}$}{
    \If{$|\text{built}| \ge B_{\text{turn}}$ \textbf{or} $\text{unstaffed}+|\text{built}| \ge U_{\max}$}{\textbf{break}}
    $s \leftarrow$ best-ranked unused site offering \text{btype}\;
    \If{$s$ exists \textbf{and} $\mathrm{cost}(s) \le \text{budget}-\beta$}{
      append build$(s)$ to $X$;\quad $\text{budget}\mathrel{-}=\mathrm{cost}(s)$;\quad mark $s$ built\;
    }
  }
}
\BlankLine
$\text{gap} \leftarrow \max(0,\ \text{need}_{\text{now}} + \omega|\text{built}| - f)$\;
\If{$\text{gap} > 0$}{
  append the largest-quantity \unity{hire\_untrained} with $\mathrm{cost}\le\text{budget}-\beta$ to $X$\;
}
\Return $(C, X)$\;
\caption{\textsc{ImprovedRulesBasedPolicy} (\py{improved\_rules\_based\_decision})}
\label{alg:v2}
\end{algorithm}

\paragraph{Empirical note.} On an 8-episode smoke (18 rounds), \py{rules-based-v2} scored
$\approx \textcolor{okgreen}{3.01}$ vs the baseline \py{rules-based} $2.49$ and the best LLM
(gpt-5.4) $2.25$: full food ($1.0$), lodging $0.93$ ($93\%$ fulfilment vs the baseline's $53\%$),
casework $0.43$ (vs $0.32$), at equal cost-efficiency and a healthier final budget. Numbers are a
smoke, not a final benchmark.

\end{document}
